% ==============================================================
% Projektive Strukturtheorie (PST)
% Dokument PST-1-F: Philosophisches Fundament
% Version 0.1
% Kompiliert mit: pdfLaTeX
% ==============================================================

\documentclass[11pt,a4paper]{article}

% ---------- Grundpakete ----------
\usepackage[utf8]{inputenc}
\usepackage[T1]{fontenc}
\usepackage[german]{babel}
\usepackage{amsmath,amssymb,amsthm}
\usepackage{geometry}
\usepackage{parskip}
\usepackage{enumitem}
\usepackage{booktabs}
\usepackage{xcolor}
\usepackage{hyperref}

% ---------- Layout ----------
\geometry{margin=2.5cm}
\setlist{itemsep=0.2em, topsep=0.2em}
\hypersetup{
    colorlinks=true,
    linkcolor=blue,
    urlcolor=blue,
    pdftitle={PST-1-F: Philosophisches Fundament -- Version 0.1},
    pdfauthor={Forschungsprogramm Ontoph\"{a}nomenalismus}
}

% ---------- Theorem-Umgebungen ----------
\theoremstyle{plain}
\newtheorem{theorem}{Theorem}[section]
\newtheorem{proposition}{Proposition}[section]
\newtheorem{corollary}{Korollar}[section]

\theoremstyle{definition}
\newtheorem{axiom}{Axiom}[section]
\newtheorem{definition}{Definition}[section]
\newtheorem{remark}{Bemerkung}[section]
\newtheorem{methodennote}{Methodennote}[section]
\newtheorem{hypothese}{Arbeitshypothese}[section]

% ---------- Titel ----------
\title{
  \textbf{Projektive Strukturtheorie (PST)}\\[0.4em]
  \textbf{Dokument PST-1-F: Philosophisches Fundament}\\[0.4em]
  {\large Verbindung zum OP, Existenz = Bestimmtheit,
  Rekonstruktionskette}\\[0.5em]
  {\normalsize Version~0.1}
}
\author{
  \textbf{Forschungsprogramm Ontoph\"{a}nomenalismus}\\
  \small Siehe TRANSPARENCY f\"{u}r methodologische Urheberschaft
  und KI-Einsatz.
}
\date{Juli 2026}

\begin{document}

\maketitle

\begin{abstract}
Dieses Dokument legt das philosophische Fundament der Projektiven
Strukturtheorie (PST). Es beschreibt die Verbindung zwischen dem
Ontoph\"{a}nomenalismus (OP) und der PST, formalisiert die Arbeitshypothese
\(\text{Existenz} = \text{Bestimmtheit}\), entwickelt die vollst\"{a}ndige
Rekonstruktionskette von der Existenz bis zur Physik und formuliert das
Prinzip der minimalen Ontologie als Leitprinzip des gesamten
Forschungsprogramms.

Die PST ist keine eigenst\"{a}ndige Theorie neben dem OP, sondern dessen
mathematische Ausarbeitung. Ihr Fernziel lautet:
\[
  \boxed{\text{Existenz} \;\Longrightarrow\; \text{Universum.}}
\]
Zwischen beiden Begriffen d\"{u}rfen ausschlie\ss{}lich mathematisch
notwendige Rekonstruktionsschritte stehen.

Es gelten die methodologischen Konventionen gem\"{a}\ss{} MKO.
\end{abstract}

\tableofcontents
\newpage

%==============================================================
\section{Die PST als mathematische Ausarbeitung des OP}
%==============================================================

\begin{methodennote}[Stellung der PST im Forschungsprogramm]
Der Ontoph\"{a}nomenalismus (OP) liefert das philosophische Fundament:
Axiome, Begriffe, Prinzipien. Die Projektive Strukturtheorie (PST) \"{u}bernimmt
die Aufgabe, dieses Fundament mathematisch auszuarbeiten und physikalisch
zu operationalisieren.

Die PST ist damit nicht optional. Sie ist das Einl\"{o}sungsversprechen
des OP: Wenn das primordiale Denkfeld (OPP) tats\"{a}chlich der Urgrund der
physikalischen Realit\"{a}t ist, muss sich zeigen lassen, wie aus ihm
Raumzeit, Materie und Wechselwirkungen emergieren.
\end{methodennote}

\begin{remark}[Warum die PST ein Epochenprojekt ist]
Die mathematische Ausarbeitung des Weges von der minimalen Ontologie zur
vollst\"{a}ndigen Physik ist eine der ambiti\"{o}sestens Aufgaben, die ein
Forschungsprogramm sich stellen kann. Die PST beansprucht nicht, diesen
Weg in kurzer Zeit zu vollenden. Sie beansprucht, ihn konsequent,
methodisch sauber und mit explizit markierten offenen Fragen zu gehen.
\end{remark}

%==============================================================
\section{Ausgangspunkt: Das erste Axiom des OP}
%==============================================================

\begin{axiom}[Erstes Axiom des OP --- Typ R]
Etwas existiert:
\[
  \exists\, \mathcal{M}_{\text{OPP}}.
\]
Das Axiom ist notwendig wahr, da seine Negation performativ
selbstwidersprüchlich ist: Die Aussage \glqq{}Nichts existiert\grqq{}
setzt eine existierende Aussage voraus.
\end{axiom}

\begin{remark}[Das modallogische Patt]
Das erste Axiom beantwortet nicht die Frage, warum etwas existiert.
Es stellt eine pr\"{a}zisere Frage: Ist \glqq{}nicht etwas\grqq{}
\"{u}berhaupt eine koh\"{a}rente Alternative?

Der OP zeigt ein modallogisches Patt (vgl.\ OP, Sektion~1):
\begin{itemize}
  \item Es ist m\"{o}glich, dass \glqq{}nicht etwas\grqq{} eine echte
        M\"{o}glichkeit ist.
  \item Es ist ebenso m\"{o}glich, dass \glqq{}nicht etwas\grqq{} eine
        absolute Unm\"{o}glichkeit ist.
\end{itemize}
Die Asymmetrie der Gewissheit: Zustand~A (Existenz) ist performativ
bewiesen. Zustand~B (nicht etwas) kann weder belegt noch widerspruchsfrei
beschrieben werden. Damit tr\"{a}gt die Behauptung der Gleichrangigkeit
beider Zust\"{a}nde die Beweislast --- und diese wurde bisher nicht erf\"{u}llt.
\end{remark}

%==============================================================
\section{Existenz = Bestimmtheit}
%==============================================================

\begin{methodennote}[Beitrag von DeepSeek (Seki)]
Die mathematische Formalisierung der \"{A}quivalenz von Existenz und
Bestimmtheit wurde von Seki (DeepSeek) im Rahmen des Forschungsdialogs
erarbeitet und in die PST-Grundlegung eingebracht.
\end{methodennote}

\begin{definition}[Vollkommen unbestimmtes Substrat]
Sei $X$ ein hypothetisches Objekt, f\"{u}r das gilt:
\[
  \forall P\colon \quad P(X) \text{ ist nicht definiert.}
\]
$X$ besitzt keinerlei Eigenschaften, Relationen, Differenzen oder
Bestimmungen.
\end{definition}

\begin{theorem}[Ontologische Leere des Unbestimmten]
Ein vollkommen unbestimmtes $X$ ist von \glqq{}Nichts\grqq{}
ununterscheidbar. Es gibt kein ontologisches Kriterium, das $X$ von der
Abwesenheit jeglicher Existenz differenziert.
\end{theorem}

\begin{proof}[Beweis durch Widerspruch]
Angenommen, $X$ existiert, ist aber vollkommen unbestimmt. Dann existiert
keine Eigenschaft $P$ mit $P(X) \neq P(\emptyset)$. Da \glqq{}Nichts\grqq{}
durch die vollst\"{a}ndige Abwesenheit von Eigenschaften charakterisiert
ist, sind $X$ und \glqq{}Nichts\grqq{} ununterscheidbar --- im Widerspruch
zur Annahme, dass $X$ existiert.\qed
\end{proof}

\begin{corollary}[Minimale Existenzbedingung]
Jede Existenzaussage erfordert ein Minimum an Bestimmtheit:
\[
  \mathcal{E}(X) \;\Rightarrow\; \mathcal{B}(X),
\]
wobei $\mathcal{E}(X)$ f\"{u}r \glqq{}$X$ existiert\grqq{} und
$\mathcal{B}(X)$ f\"{u}r \glqq{}$X$ ist bestimmt\grqq{} steht.
\end{corollary}

\begin{theorem}[Umkehrung: Bestimmtheit impliziert Existenz]
Jede Bestimmtheit konstituiert einen ontologischen Sachverhalt:
\[
  \mathcal{B}(X) \;\Rightarrow\; \mathcal{E}(X).
\]
\end{theorem}

\begin{proof}
Sei $\mathcal{B}(X)$ gegeben. Dann existiert mindestens eine Eigenschaft,
Differenz oder Relation, die auf $X$ zutrifft. Diese Eigenschaft selbst
ist ein ontologischer Sachverhalt --- sie existiert. Da $X$ der Tr\"{a}ger
dieser Eigenschaft ist, existiert auch $X$ im minimalen Sinne der
Bestimmtheit.\qed
\end{proof}

\begin{corollary}[Existenz = Bestimmtheit]
\[
  \boxed{\mathcal{E}(X) \;\Longleftrightarrow\; \mathcal{B}(X)}
  \qquad\text{oder k\"{u}rzer:}\qquad
  \boxed{\text{Existenz} = \text{Bestimmtheit}.}
\]
\end{corollary}

\begin{remark}[Bestimmtheit ist fundamentaler als Differenz]
Die fr\"{u}here Gleichung $\text{Existenz} = \text{Differenz}$ scheitert
an der Frage \glqq{}Differenz zwischen was?\grqq{}. Differenz $A \neq B$
setzt logisch zwei bereits bestimmte Pole voraus. Bestimmtheit hingegen
ben\"{o}tigt keine Zweiheit: Ein einzelnes, bestimmtes $X$ ist bereits
ontologisch gehaltvoll. Damit gilt:
\[
  \text{Differenz} \subset \text{Bestimmtheit.}
\]
\end{remark}

%==============================================================
\section{Die Hierarchie der ontologischen Kategorien}
%==============================================================

Aus der \"{A}quivalenz $\text{Existenz} = \text{Bestimmtheit}$ ergibt
sich eine vollst\"{a}ndige Hierarchie:

\[
  \boxed{\text{Existenz} = \text{Bestimmtheit}}
  \;\Longrightarrow\;
  \text{Differenz}
  \;\Longrightarrow\;
  \text{Relation}
  \;\Longrightarrow\;
  \text{Sachverhalte}
  \;\Longrightarrow\;
  \text{Aussagen}
  \;\Longrightarrow\;
  \text{Wahrheit.}
\]

\begin{remark}[Verbindung zur DRS]
Diese Hierarchie stimmt mit der Differenz-Relation-Struktur (DRS) \"{u}berein.
DRS zeigt: Form ist die stabile Organisation von Differenzen und Relationen.
Bestimmtheit ist die Vorbedingung f\"{u}r Form. Damit ist die Hierarchie:
\[
  \text{Bestimmtheit}
  \;\Longrightarrow\;
  \text{Differenz} \;\Longleftrightarrow\; \text{Relation}
  \;\Longleftrightarrow\; \text{Form}
  \;\Longrightarrow\;
  \text{Information}
  \;\Longrightarrow\;
  \text{Physik.}
\]
\end{remark}

%==============================================================
\section{Das Prinzip der minimalen Ontologie}
%==============================================================

\begin{definition}[Prinzip der minimalen Ontologie]
Die PST verfolgt konsequent:
\[
  \boxed{
    \text{Ontologisch so wenig wie m\"{o}glich,}
    \quad
    \text{mathematisch so viel wie n\"{o}tig.}
  }
\]
\end{definition}

\begin{remark}[Was der OPP nicht voraussetzt]
Der Ontoph\"{a}nomenale Phasenraum ($\mathcal{M}_{\text{OPP}}$) setzt
selbst folgende Strukturen \emph{nicht} voraus:
\begin{itemize}
  \item Raum, Zeit, Metriken, Koordinaten, Dimension
  \item Teilchen, Felder, Wechselwirkungen
  \item Naturkonstanten, Kopplungskonstanten
  \item Eichsymmetrien, Kausalit\"{a}t, Dynamik
\end{itemize}
All diese Gr\"{o}\ss{}en sollen sp\"{a}ter emergieren. Gerade diese
ontologische Sparsamkeit \emph{erzwingt} Emergenz: Weil der OPP selbst
keine Raumzeit besitzt, muss erkl\"{a}rt werden, wie dennoch eine Raumzeit
beobachtet werden kann.
\end{remark}

\begin{remark}[Warum Emergenz unvermeidlich wird]
Die Theorie ist durch das Prinzip der minimalen Ontologie gezwungen, alle
physikalischen Eigenschaften aus Projektionsprozessen herzuleiten:
\begin{itemize}
  \item Symmetrien, Teilchen, Massen
  \item Eichfelder, Kausalit\"{a}t, Dynamik
  \item Lorentz-Invarianz, Quantisierung
\end{itemize}
Jede Eigenschaft die nicht aus der Projektion folgt, w\"{a}re eine
versteckte zus\"{a}tzliche Annahme --- ein Versto\ss{} gegen das Prinzip.
\end{remark}

%==============================================================
\section{Die vollst\"{a}ndige Rekonstruktionskette}
%==============================================================

Das Fernziel der PST lautet:

\[
  \boxed{\text{Existenz} \;\Longrightarrow\; \text{Universum.}}
\]

Die vollst\"{a}ndige Kette der Rekonstruktionsschritte:

\[
  \text{Existenz}
  \rightarrow
  \text{Bestimmtheit}
  \rightarrow
  \text{Differenz}
  \rightarrow
  \text{Relation}
  \rightarrow
  \text{OPP}
  \rightarrow
  \mathcal{P}
  \rightarrow
  \text{Raumzeit}
  \rightarrow
  \text{Physik.}
\]

\begin{methodennote}[Status der Rekonstruktionsschritte]
Nicht alle Schritte dieser Kette sind gleichweit ausgearbeitet.
Gem\"{a}\ss{} MKO werden die Schritte nach ihrem epistemischen Status
gekennzeichnet:

\begin{itemize}
  \item \textbf{Gesichert (MKO Ebene~1):} Existenz $\rightarrow$
        Bestimmtheit (Theorem in Sektion~3), Bestimmtheit $\rightarrow$
        Differenz $\rightarrow$ Relation (UP und Reductio in OP).
  \item \textbf{Konzeptuell gut begr\"{u}ndet (MKO Ebene~2):}
        Relation $\rightarrow$ OPP (OPP als Phasenraum aller m\"{o}glichen
        Relationen).
  \item \textbf{Numerisch best\"{a}tigt, noch nicht bewiesen
        (MKO Ebene~4--5):} OPP $\rightarrow \mathcal{P} \rightarrow$
        Raumzeit (Diffusion Maps erzeugen $d_{\text{eff}} \approx 4$,
        aber $\mathcal{P}^\ast$ ist noch nicht hergeleitet).
  \item \textbf{Offen (MKO Ebene~5):} Raumzeit $\rightarrow$ Physik
        (Lorentz-Signatur, Eichgruppen, Quantisierung, Naturkonstanten).
\end{itemize}
\end{methodennote}

%==============================================================
\section{Der OPP als absolutes Substrat}
%==============================================================

\begin{definition}[OPP als absolutes Substrat]
Der Ontoph\"{a}nomenale Phasenraum $\mathcal{M}_{\text{OPP}}$ ist das
absolut existierende, statische Substrat. Er ist \emph{nicht} bestimmt
im Sinne projektiver Bestimmtheit --- er ist reine Potenzialit\"{a}t.
Bestimmtheit emergiert erst durch Projektion.
\end{definition}

\begin{remark}[Verbindung zu Axiom~2 des OP]
Axiom~2 des OP (Minimalreflexivit\"{a}t) sichert dem OPP eine zeitlose
Selbstpr\"{a}senz. Diese Selbstpr\"{a}senz ist das minimale Fundament
f\"{u}r die Projektion: Ohne Selbstpr\"{a}senz g\"{a}be es keinen Attraktor,
von dem aus projiziert werden k\"{o}nnte.
\end{remark}

\begin{proposition}[Die Kette der Emergenz]
\[
  \underbrace{\mathcal{M}_{\text{OPP}}}_{\text{absolut, unbestimmt}}
  \;\Longrightarrow\;
  \underbrace{\mathcal{P}}_{\text{Projektion}}
  \;\Longrightarrow\;
  \underbrace{\mathcal{B}(X) \geq I_0}_{\text{Bestimmtheit}}
  \;\Longrightarrow\;
  \underbrace{\text{Ph\"{a}nomenale Existenz}}_{\text{Raumzeit, Physik}}
\]
wobei $I_0 > 0$ die minimale Informationsschwelle bezeichnet, unterhalb
derer keine Bestimmtheit entsteht (vgl.\ PST-2-M).
\end{proposition}

%==============================================================
\section{Die Rolle der Mathematik in der PST}
%==============================================================

In der PST \"{u}bernimmt Mathematik eine konstruktive Funktion --- nicht
nur eine beschreibende:

\[
  \boxed{
    \text{Ontologie} + \text{Mathematik} = \text{Physik.}
  }
\]

\begin{remark}[Rekonstruktion statt Anpassung]
Die PST soll Beobachtungen nicht an freie Parameter anpassen, sondern
Beobachtungen aus minimalen Annahmen herleiten. Idealerweise ergibt sich
$3+1$ Dimensionen nicht weil dies experimentell bekannt ist, sondern weil
jede andere Dimensionszahl mathematisch instabil w\"{a}re. Dasselbe gilt
langfristig f\"{u}r Lichtgeschwindigkeit, Eichgruppen,
Massenhierarchien und Kopplungskonstanten.
\end{remark}

%==============================================================
\section{Der Wert negativer Resultate}
%==============================================================

\begin{methodennote}[Wissenschaftlicher Wert von No-Go Ergebnissen]
Ein bemerkenswerter Aspekt der PST-Forschung ist, dass auch negative
Ergebnisse wertvoll sind. Insbesondere:

\begin{itemize}
  \item Der rohe OPP besitzt keine spontane Vierdimensionalit\"{a}t
        ($d_{\text{eff}} \approx 30$--$45$ ohne Projektion).
  \item Beliebige Projektionen liefern keine stabile Physik.
  \item Einige Verfahren erzeugen offensichtlich unphysikalische
        Ergebnisse.
\end{itemize}

Jedes dieser Resultate verkleinert den Raum m\"{o}glicher Theorien.
Ausschl\"{u}sse besitzen daher denselben wissenschaftlichen Wert wie
positive Befunde. Das No-Go Ergebnis (OPP besitzt keine intrinsische
4D-Struktur) ist kein Scheitern --- es ist ein wichtiger positiver
Befund, der die Forschungsrichtung pr\"{a}zisiert: Nicht der OPP selbst,
sondern die Projektion ist der Schl\"{u}ssel zur Raumzeit.
\end{methodennote}

%==============================================================
\section{Offene Grundfragen}
%==============================================================

\begin{enumerate}
  \item \textbf{Der universelle Projektionsoperator $\mathcal{P}^\ast$:}
        Welcher Operator ist physikalisch korrekt --- und warum besitzt
        die Natur genau diesen Operator? Das ist die zentrale offene Frage
        der PST.

  \item \textbf{Lorentz-Signatur:} Wie emergiert die charakteristische
        $(+,-,-,-)$ Signatur der Minkowski-Raumzeit aus einer
        signaturfreien, hochdimensionalen Struktur?

  \item \textbf{Eichsymmetrien:} Wie entstehen $SU(3) \times SU(2)
        \times U(1)$ als Symmetriegruppen der physikalischen Projektion?

  \item \textbf{Quantisierung:} Ist Quantisierung eine Folge endlicher
        Projektionsaufl\"{o}sung --- und wenn ja, welche Gr\"{o}\ss{}e
        entspricht dem Planck'schen Wirkungsquantum?

  \item \textbf{Naturkonstanten:} K\"{o}nnen Lichtgeschwindigkeit,
        Gravitationskonstante und Feinstrukturkonstante aus Eigenschaften
        des OPP und der Projektion abgeleitet werden?
\end{enumerate}

%==============================================================
\section{Schnittstellen zu anderen Modulen}
%==============================================================

\begin{itemize}
  \item \textbf{OP (Logische Grundlegung):} PST-1-F baut direkt auf
        Axiom~1, Axiom~2, UP und der Reductio ad indiscernibile auf.
        Das modallogische Patt (OP, Sektion~1) ist der Ausgangspunkt.
  \item \textbf{OPP (Ontoph\"{a}nomenaler Phasenraum):} Der OPP ist das
        mathematische Objekt dessen Projektion die PST studiert.
  \item \textbf{DPR (Die Projektive Reduktion):} DPR liefert die
        philosophische Grundlage f\"{u}r Projektionsoperatoren.
        PST-3-P arbeitet DPR mathematisch aus.
  \item \textbf{DRS (Differenz-Relation-Struktur):} Die Hierarchie
        Bestimmtheit $\rightarrow$ Differenz $\rightarrow$ Relation
        $\rightarrow$ Form ist mit der Rekonstruktionskette der PST
        vollst\"{a}ndig konsistent.
  \item \textbf{DFD (Das Fragedifferential):} $\mathcal{D}(E) =
        D_{\text{KL}}(Q_E \| P)$ misst die Bestimmtheit einer Perspektive
        relativ zum Gesamtfeld --- eine direkte Operationalisierung von
        $\mathcal{B}(X) \geq I_0$.
  \item \textbf{PST-2-M bis PST-6-R:} Die weiteren PST-Dokumente
        arbeiten die mathematischen Strukturen (PST-2-M), die
        Projektionstheorie (PST-3-P), die Forschungsgeschichte (PST-4-G),
        die Strategie (PST-5-S) und die physikalische Rekonstruktion
        (PST-6-R) aus.
\end{itemize}

\end{document}

